% Shared preamble of the Mass booklets: % Shared preamble of the Mass booklets: % Shared preamble of the Mass booklets: % Shared preamble of the Mass booklets: \input{../shared/booklet} first.
\documentclass[a4paper,11pt]{article}

\usepackage[a4paper,top=15mm,bottom=18mm,left=20mm,right=20mm]{geometry}
\usepackage{fontspec}
\setmainfont{Libertinus Serif}
\usepackage[latin,main=polish]{babel}
\usepackage{microtype}
\usepackage{xcolor}
\usepackage{paracol}
\usepackage{titlesec}
\usepackage[
  program=lilypond,
  staffsize=16,
  includepaths={./,scores/,../shared/,../shared/scores/},
  pass-fonts,
  tmpdir=build-ly,
]{lyluatex}

\setlength{\parindent}{0pt}
\setlength{\parskip}{0pt}
\pagestyle{plain}

% \chant{Latin title}{mode}
\titleformat{\section}{\centering\Large\scshape}{}{0pt}{}
\titlespacing*{\section}{0pt}{1.8ex plus 1ex}{0.3ex}
\newcommand{\chant}[2]{%
  \section{\foreignlanguage{latin}{#1}}%
  {\centering\small\itshape #2\par}%
  \vspace{0.8ex}}

\newcommand{\tekstsize}{\small}
\input{../shared/macros}

% Credo and the Marian antiphons heading are printed in place in a booklet
% (the Kyriale prints them in their own part).
\newcommand{\credo}[1]{\clearpage\input{../shared/credo-#1}}
\newcommand{\antiphonsheading}{%
  \clearpage
  \begin{center}
    {\LARGE\scshape Antyfony maryjne}\par\smallskip
    {\small\itshape śpiewane po Mszy świętej, według okresu roku liturgicznego}
  \end{center}}
 first.
\documentclass[a4paper,11pt]{article}

\usepackage[a4paper,top=15mm,bottom=18mm,left=20mm,right=20mm]{geometry}
\usepackage{fontspec}
\setmainfont{Libertinus Serif}
\usepackage[latin,main=polish]{babel}
\usepackage{microtype}
\usepackage{xcolor}
\usepackage{paracol}
\usepackage{titlesec}
\usepackage[
  program=lilypond,
  staffsize=16,
  includepaths={./,scores/,../shared/,../shared/scores/},
  pass-fonts,
  tmpdir=build-ly,
]{lyluatex}

\setlength{\parindent}{0pt}
\setlength{\parskip}{0pt}
\pagestyle{plain}

% \chant{Latin title}{mode}
\titleformat{\section}{\centering\Large\scshape}{}{0pt}{}
\titlespacing*{\section}{0pt}{1.8ex plus 1ex}{0.3ex}
\newcommand{\chant}[2]{%
  \section{\foreignlanguage{latin}{#1}}%
  {\centering\small\itshape #2\par}%
  \vspace{0.8ex}}

\newcommand{\tekstsize}{\small}
% Macros shared by the Mass booklets and the Kyriale; the including preamble
% loads fontspec, babel, xcolor and paracol and defines \tekstsize.
\definecolor{rubric}{RGB}{170,20,20}

% Parallel text: Latin left, Polish right, one row per sung phrase.
\newenvironment{tekst}
  {\par\vspace{0.6ex}\tekstsize
   \setcolumnwidth{0.48\textwidth,0.48\textwidth}%
   \setlength{\columnsep}{0.04\textwidth}%
   \begin{paracol}{2}}
  {\end{paracol}\par\vspace{0.5ex}}
% \w{Latin}{Polski}
\newcommand{\w}[2]{%
  \switchcolumn[0]*\foreignlanguage{latin}{#1}\par
  \switchcolumn #2\par\vspace{0.3ex}}
% Libertinus has no U+211F/U+2123: build the barred R and V.
\newcommand{\barred}[1]{#1\kern-.48em\raisebox{.04em}{/}\kern.06em}
\newcommand{\R}{\textcolor{rubric}{\barred{R}.}~}
\newcommand{\V}{\textcolor{rubric}{\barred{V}.}~}
\newcommand{\rubr}[1]{\textcolor{rubric}{\itshape #1}}


% Credo and the Marian antiphons heading are printed in place in a booklet
% (the Kyriale prints them in their own part).
\newcommand{\credo}[1]{\clearpage\input{../shared/credo-#1}}
\newcommand{\antiphonsheading}{%
  \clearpage
  \begin{center}
    {\LARGE\scshape Antyfony maryjne}\par\smallskip
    {\small\itshape śpiewane po Mszy świętej, według okresu roku liturgicznego}
  \end{center}}
 first.
\documentclass[a4paper,11pt]{article}

\usepackage[a4paper,top=15mm,bottom=18mm,left=20mm,right=20mm]{geometry}
\usepackage{fontspec}
\setmainfont{Libertinus Serif}
\usepackage[latin,main=polish]{babel}
\usepackage{microtype}
\usepackage{xcolor}
\usepackage{paracol}
\usepackage{titlesec}
\usepackage[
  program=lilypond,
  staffsize=16,
  includepaths={./,scores/,../shared/,../shared/scores/},
  pass-fonts,
  tmpdir=build-ly,
]{lyluatex}

\setlength{\parindent}{0pt}
\setlength{\parskip}{0pt}
\pagestyle{plain}

% \chant{Latin title}{mode}
\titleformat{\section}{\centering\Large\scshape}{}{0pt}{}
\titlespacing*{\section}{0pt}{1.8ex plus 1ex}{0.3ex}
\newcommand{\chant}[2]{%
  \section{\foreignlanguage{latin}{#1}}%
  {\centering\small\itshape #2\par}%
  \vspace{0.8ex}}

\newcommand{\tekstsize}{\small}
% Macros shared by the Mass booklets and the Kyriale; the including preamble
% loads fontspec, babel, xcolor and paracol and defines \tekstsize.
\definecolor{rubric}{RGB}{170,20,20}

% Parallel text: Latin left, Polish right, one row per sung phrase.
\newenvironment{tekst}
  {\par\vspace{0.6ex}\tekstsize
   \setcolumnwidth{0.48\textwidth,0.48\textwidth}%
   \setlength{\columnsep}{0.04\textwidth}%
   \begin{paracol}{2}}
  {\end{paracol}\par\vspace{0.5ex}}
% \w{Latin}{Polski}
\newcommand{\w}[2]{%
  \switchcolumn[0]*\foreignlanguage{latin}{#1}\par
  \switchcolumn #2\par\vspace{0.3ex}}
% Libertinus has no U+211F/U+2123: build the barred R and V.
\newcommand{\barred}[1]{#1\kern-.48em\raisebox{.04em}{/}\kern.06em}
\newcommand{\R}{\textcolor{rubric}{\barred{R}.}~}
\newcommand{\V}{\textcolor{rubric}{\barred{V}.}~}
\newcommand{\rubr}[1]{\textcolor{rubric}{\itshape #1}}


% Credo and the Marian antiphons heading are printed in place in a booklet
% (the Kyriale prints them in their own part).
\newcommand{\credo}[1]{\clearpage\input{../shared/credo-#1}}
\newcommand{\antiphonsheading}{%
  \clearpage
  \begin{center}
    {\LARGE\scshape Antyfony maryjne}\par\smallskip
    {\small\itshape śpiewane po Mszy świętej, według okresu roku liturgicznego}
  \end{center}}
 first.
\documentclass[a4paper,11pt]{article}

\usepackage[a4paper,top=15mm,bottom=18mm,left=20mm,right=20mm]{geometry}
\usepackage{fontspec}
\setmainfont{Libertinus Serif}
\usepackage[latin,main=polish]{babel}
\usepackage{microtype}
\usepackage{xcolor}
\usepackage{paracol}
\usepackage{titlesec}
\usepackage[
  program=lilypond,
  staffsize=16,
  includepaths={./,scores/,../shared/,../shared/scores/},
  pass-fonts,
  tmpdir=build-ly,
]{lyluatex}

\setlength{\parindent}{0pt}
\setlength{\parskip}{0pt}
\pagestyle{plain}

% \chant{Latin title}{mode}
\titleformat{\section}{\centering\Large\scshape}{}{0pt}{}
\titlespacing*{\section}{0pt}{1.8ex plus 1ex}{0.3ex}
\newcommand{\chant}[2]{%
  \section{\foreignlanguage{latin}{#1}}%
  {\centering\small\itshape #2\par}%
  \vspace{0.8ex}}

\newcommand{\tekstsize}{\small}
% Macros shared by the Mass booklets and the Kyriale; the including preamble
% loads fontspec, babel, xcolor and paracol and defines \tekstsize.
\definecolor{rubric}{RGB}{170,20,20}

% Parallel text: Latin left, Polish right, one row per sung phrase.
\newenvironment{tekst}
  {\par\vspace{0.6ex}\tekstsize
   \setcolumnwidth{0.48\textwidth,0.48\textwidth}%
   \setlength{\columnsep}{0.04\textwidth}%
   \begin{paracol}{2}}
  {\end{paracol}\par\vspace{0.5ex}}
% \w{Latin}{Polski}
\newcommand{\w}[2]{%
  \switchcolumn[0]*\foreignlanguage{latin}{#1}\par
  \switchcolumn #2\par\vspace{0.3ex}}
% Libertinus has no U+211F/U+2123: build the barred R and V.
\newcommand{\barred}[1]{#1\kern-.48em\raisebox{.04em}{/}\kern.06em}
\newcommand{\R}{\textcolor{rubric}{\barred{R}.}~}
\newcommand{\V}{\textcolor{rubric}{\barred{V}.}~}
\newcommand{\rubr}[1]{\textcolor{rubric}{\itshape #1}}


% Credo and the Marian antiphons heading are printed in place in a booklet
% (the Kyriale prints them in their own part).
\newcommand{\credo}[1]{\clearpage\input{../shared/credo-#1}}
\newcommand{\antiphonsheading}{%
  \clearpage
  \begin{center}
    {\LARGE\scshape Antyfony maryjne}\par\smallskip
    {\small\itshape śpiewane po Mszy świętej, według okresu roku liturgicznego}
  \end{center}}
